\documentclass{article}
\usepackage[utf8]{inputenc}

\title{Práctica 1 \\ Introducción a las herramientas del laboratorio}
\date{Fecha limite de entrega 15 de agosto a las 23:59}
\usepackage{natbib}
\usepackage{graphicx}
\providecommand{\tightlist}{}

\begin{document}

\maketitle

\section{Objetivo}
{Debido a que vamos a usar este repositorio para absolutamente todo, esta práctica es }{OBLIGATORIA}\\\\

El objetivo de esta práctica es familiarizarse con las tecnologías que se usarán en el 100\% de las siguientes prácticas, así como familiarizarse con el formato de las entregas de las prácticas y los proyectos. \\\\
Al terminar la práctica el alumno tendrá un repositorio personal para el curso de ICC y tendrá una noción general del proceso de programación.

\section{Desarrollo}
Antes de comenzar a programar debemos comenzar a aprender a usar el ambiente de desarrollo, y las herramientas necesarias para esta tarea.

{Las principales herramientas serán:}

\begin{itemize}
\tightlist
\item
  {La consola de tu sistema operativo}
\item
  {Git}
\item
  {El compilador de Java}
\item
  {El intérprete de la JVM}
\item
  {Editor de texto de tu preferencia}
\end{itemize}


{Primero lo primero}

\begin{enumerate}
\tightlist
\item
  {El primer paso de la práctica es crear repositorio público en github con el siguiente nombre }{obligatoriamente}
\end{enumerate}

{ICC\_20201\_NUMCUENTA }{donde NUMCUENTA es tu número de cuenta.}

{ejemplo* ~ ~ ~ ~ ~ ~ ~ICC\_20201\_301233514}
\\
{El repositorio será público, y no debe ser inicializado con ningún README}

{}

{}

{}

{}

{}

\begin{enumerate}
\setcounter{enumi}{1}
\tightlist
\item
  {Crear un documento en LaTex de tipo reporte contestando las siguientes preguntas}
\end{enumerate}

\begin{enumerate}
\tightlist
\item
  {¿Qué es una clase?}
\item
  {¿Cúal es la diferencia entre una clase y un objeto?}
\item
  {Dar 3 ejemplos de clases y un par de objetos por cada clase}
\item
  {¿Cúal es el propósito de la documentación?}
\item
  {¿Qué ventajas tiene hacer el diseño de un programa antes de codificar?}
\item
  {¿Qué comando se usa para ejecutar un programa en java?}
\item
  {¿Qué comando se usa para compilar un programa en java?}
\end{enumerate}

{}

\begin{enumerate}
\setcounter{enumi}{2}
\tightlist
\item
  {En la carpeta que hayas elegido para todas tus prácticas y proyectos crea un archivo README.txt, que tendrá tu número de cuenta y tu nombre, dentro de esa carpeta, crea otra carpeta que se llame practica01 y dentro de practica01, ubica el pdf del paso 2}
\end{enumerate}

{}

{}

{Las carpetas estarán de la siguiente manera.}

\begin{verbatim}
./
|-- practica01
|   `-- preguntas.pdf
`-- README.txt
\end{verbatim}

{}

\begin{enumerate}
\setcounter{enumi}{3}
\tightlist
\item
  {Subir toda la carpeta anterior al repositorio creado en github}
\item
  {U}{na vez que tengas tu repositorio listo, copiar la dirección del repositorio y enviarla por correo a la dirección iabin@ciencias.unam.mx}
\end{enumerate}


\end{document}
